%%%%%%%%%%%%%%%%%%%%%%%%%%%%%%%%%%%%%%%%%%%%%%%%%%%%%%%%%%%%%%%%%%%%%%%%%%%%%%%%%%%%%%%%%%%%%%%%%%%%%%
%
%   Filename    : chapter_1.tex 
%
%   Description : This file will contain your Research Description.
%                 
%%%%%%%%%%%%%%%%%%%%%%%%%%%%%%%%%%%%%%%%%%%%%%%%%%%%%%%%%%%%%%%%%%%%%%%%%%%%%%%%%%%%%%%%%%%%%%%%%%%%%%

\chapter{Introduction}
\label{sec:intro}    %--note: labels help you with hyperlink editing (using your IDE)


%%
%% --- 1.1 Background of the Study --- %%
%%

\section{Background of the Study}
\label{sec:overview}

%
%   NOTE: You have to delete/replace the unnecessary paragraphs with your own text.
%

% Describe the research area

Volcanic eruptions are among the most destructive natural hazards, producing ash plumes, pyroclastic flows, and volcanic gases that threaten nearby communities and critical infrastructure. 
Effective visualization of these hazards plays an important role in disaster preparedness, scientific communication, and public education.
Traditionally, volcanic simulations have relied on numerical models that prioritize scientific accuracy for hazard prediction and analysis.
While these models provide valuable quantitative information, they are often computationally intensive and lack the interactive capabilities necessary for real-time visualization \cite{rossant2021high}.

% Describe the domain (where will you apply the topic)

To overcome these limitations, the subfield of real-time computer graphics has increasingly been applied to scientific visualization, balancing visual plausibility with high interactive performance.
Within this domain, a notable localized development is AnitoPlume, a real-time three-dimensional volcanic plume simulator designed specifically for visualizing the complex eruption dynamics of Taal Volcano \cite{del2026interactive}.

\begin{figure}[htbp]
    \centering
    \makebox[\textwidth][c]{
        \includegraphics[width=1.2\textwidth]{AnitoPlume-Interface.png}
    }
    \caption{User Interface of AnitoPlume. The Taal volcano model is rendered in the center of the simulator. The photo reference is on the left. \protect\cite{del2026interactive}.}
    \label{fig:single_experiment}
\end{figure}

This specialized application combines a Lagrangian particle-based plume system with a detailed digital elevation model of the volcano's terrain.
By maintaining an interactive rendering performance of approximately 30 frames per second, the system demonstrated that computer graphics could bridge the gap between complex geological events and accessible hazard-awareness software for local education and planning \cite{del2026interactive}.

% Give a synthesis of previous work

Previous efforts in graphics-driven volcanic visualization have heavily relied on traditional, high-level graphics architectures due to their broad portability and structural simplicity \cite{Shiraef2016}.
In these legacy state-machine models, the graphics API simplifies software development by delegating critical hardware management tasks, such as memory allocation and command scheduling, directly to the vendor's graphics driver.
However, as visualization workloads scale, this abstraction inevitably introduces massive CPU driver overhead and restricts developers from optimizing how the hardware executes commands \cite{bossard2019high}.
Consequently, the computer graphics industry has seen a global shift toward modern, explicit rendering paradigms. 
These low-level architectures strip away hidden driver intervention and grant developers granular, direct control over GPU resources, pipeline states, and memory synchronization to maximize hardware efficiency \cite{ioannidis2020multithreaded}.

% State the problem and the research question

\begin{figure}[htbp]
    \centering
    \makebox[\textwidth][c]{
        \includegraphics[width=0.85\textwidth]{Graphics_APIs.png}
    }
    % \protect is mandatory when using apacite inside captions
    \caption{Graphics APIs. \protect\cite{TechArx2016}.}
    \label{fig:graphics_api}
\end{figure}

Despite the theoretical performance advantages of modern explicit rendering pipelines, a significant challenge remains: the architectural enhancement of established scientific tools.
High-level visualization platforms like the original AnitoPlume simulator are tightly coupled to legacy architectures, creating a rigid barrier to optimization and preventing the integration of advanced, contemporary rendering techniques. 
Transitioning an active simulation system to an explicit framework requires a complete restructuring of shader programs, command submissions, and memory allocation strategies. 
This process moves beyond a mere software translation; it presents a critical opportunity to upgrade the shading architecture and optimize the visual representation of complex fluid-terrain interactions \cite{Unterguggenberger2022}.
Therefore, investigating how explicit GPU control impacts rendering throughput, memory bandwidth utilization, and overall visual fidelity provides essential insights for the future development of highly scalable scientific visualization systems.
%%
%% --- 1.2 Research Objectives --- %%
%%

\section{Research Objectives}

This study aims to port and extend AnitoPlume, an existing Taal Volcano plume simulator, by redesigning its rendering architecture with a modern graphics framework while preserving its established particle-based simulation model. The research evaluates the resulting improvements in rendering performance, GPU resource utilization, scalability, and visual fidelity, and assesses whether the modernized system maintains an interactive and usable experience for its target beneficiaries.\begin{comment}
%
% IPR acknowledgement: the following sentences and examples are from Ethel Ong's slides 
%     on Research Objectives
%

How to formulate your research objectives:
1. Identify what research steps do you need to perform to achieve your general objective.
2. Identify the questions that must be answered for you to achieve your general objective.
    Thereafter, convert these questions into action statements


Example #1:

Question:
    What strategies do human educators employ in collaborative storytelling with children?

Specific Objective:
   To review existing strategies employed by language educators when sharing storytelling with children


Example #2:

Question:
   How will you represent commonsense knowledge for use by computer systems?

Specific Objective:
   To identify knowledge representation approaches used by existing story generation systems

Example #3:
Question:
   What types of storytelling knowledge are needed to generate stories?

Objective:
    To identify the different types of storytelling knowledge used in generating stories

Example #4:
Question:
    What machine learning approaches will you utilize?

Specific Objective:
    To determine existing machine learning algorithms [that can be used in training the computer system to detect cyberbullying cases] 

Example #5: 
Question:
    How will your research output be evaluated?

Specific Objective:
    To define evaluation metrics for validating the accuracy of the translation
\end{comment}


%
%  The following is the format for presenting your specific objectives; replace them with your own 
%

The specific objectives of this research are as follows:
\begin{enumerate}
   \item To benchmark GPU resource synchronization and cross-API shader translation for AnitoPlume's volumetric plume data, establishing a baseline for porting the simulator to the new rendering framework;
   \item To redesign AnitoPlume's software architecture by transitioning from a dynamic state-machine model to explicitly defined Pipeline State Objects (PSOs) and command lists;
   \item To analyze the rendering efficiency of the modernized pipeline against the previous implementation using GPU profiling metrics, such as, but not limited to, frame pacing, shader execution time, memory bandwidth utilization, and overall render throughput, under identical hardware constraints;
   \item To validate the overall user experience of the modernized simulator by conducting usability testing with target beneficiaries.
\end{enumerate}


%%
%% --- 1.3 Scope and Limitations --- %%
%%

\section{Scope and Limitations of the Research}
\label{sec:scopelimitations}

This study is limited to porting and extending AnitoPlume, an existing Taal Volcano plume visualization system, from its current graphics API (OpenGL) to a modern graphics API, such as Vulkan or DirectX 12. In contrast to legacy APIs, where the graphics driver automatically manages GPU state, memory allocation, and synchronization on the application's behalf, explicit APIs require the application to directly define and control these aspects, including Pipeline State Objects (PSOs), command lists, and GPU resource synchronization. This shift is the basis for the "explicit" rendering architecture investigated in this research. The study specifically examines pipeline state management, command recording, shader portability across graphics APIs, and GPU resource synchronization, while strictly preserving AnitoPlume's existing particle-based plume simulation model. Modifications to the underlying scientific simulation algorithms, volcanic plume physics, or eruption modeling parameters are outside the scope of this work.

The usability testing phase will involve 10 to 20 voluntary participants, consisting of computer science students and developers with a relevant background in 3D graphics or software development. Participants will be recruited on-campus within DLSU computer laboratories. Data collection will take place within the DLSU College of Computer Studies (CCS) laboratory environments, or in controlled remote developer testing environments for participants unable to be physically present. After completing a standardized set of interaction tasks with the ported simulator, participants will complete an anonymous, self-administered questionnaire using the System Usability Scale (SUS) to evaluate application responsiveness, stability, and visual parity relative to the original AnitoPlume system. This phase is expected to take 2 to 4 weeks within the systems evaluation timeframe of the study.

This research is limited to a single case study, the AnitoPlume Taal Volcano plume simulator, and its findings on rendering performance and usability are not guaranteed to generalize to other volcanic visualization systems or simulation platforms. Benchmarking results are further constrained to the specific hardware configurations used during testing and may not be representative of performance across all GPU vendors or consumer hardware. Finally, given the sample size of 10 to 20 participants, usability findings are intended to provide indicative, rather than statistically generalizable, insight into user experience.
%%
%% --- 1.4 Significance of the Research --- %%
%%

\section{Significance of the Research}
This study contributes to both computer graphics research and scientific visualization by providing empirical evidence on the application of modern graphics architectures to real-time volumetric scientific visualization, while also offering immediate, practical value to Philippine disaster risk reduction efforts.

For Philippine communities and disaster management stakeholders, this study offers direct, immediate benefits. Taal Volcano remains under close monitoring due to its recurring eruptive activity, yet accessible, interactive 3D tools for communicating volcanic hazards to the public and to local decision-makers remain limited. By modernizing and extending AnitoPlume into a more performant, maintainable, and scalable system, this study lays the groundwork for AnitoPlume V2 to potentially serve as one of the first accessible 3D visualization tools for volcanic hazard communication in the Philippines. Such a tool could support disaster preparedness education, hazard awareness campaigns, and public information efforts led by local government units (LGUs) and agencies such as PHIVOLCS, by allowing communities and responders to visualize eruption scenarios in a more intuitive and immersive manner than static hazard maps alone. While this study does not directly evaluate deployment with disaster response agencies or affected communities, it establishes the technical foundation necessary to pursue such applications in future work.

For researchers in computer graphics, the study offers a detailed case study demonstrating the architectural modifications required when transitioning between two fundamentally different graphics APIs.

For developers of scientific visualization software, the findings may serve as a reference for future projects that require designing high-performance rendering systems for similarly data-intensive, real-time visualization applications.

For future researchers, this work establishes a scalable rendering foundation that supports integrating advanced graphics techniques, including hardware-accelerated ray tracing, compute shaders, GPU-driven rendering, and other contemporary visualization methods. Consequently, the study extends the long-term applicability of the original AnitoPlume project beyond its initial legacy implementation, while positioning it as a potential stepping stone toward practical, real-world disaster communication tools for volcanic hazards in the Philippines.
